\section{Introduction}\label{sec:introduction56}
A continuation approximation matters because it changes an economic decision. Its fitting error, the loss of its implemented policy, and the work needed to certify an improvement are different objects. This paper develops Neural Bellman Operators around that distinction. A construction uses its own fitted future values, returns feasible action witnesses, and connects numerical resources to the cost of the policy actually implemented.

There are two linked numerical contracts in the same controlled economy. Under effective primitive moduli, the constructive backend returns a feasible policy with a prescribed error account against the original Bellman optimum. Theorems~\ref{thm:construct49} and \ref{thm:compile49} retain the construction and its exact continuous original-owner implementation. A fitted backend instead proposes actions from learned features. Those proposals become deployable only after a whole-cell policy certificate is satisfied. A prospective service then stops when the complete policy attains a prespecified expected-cost target relative to an installed controller. This additional contract neither changes the Bellman optimum nor turns an empirical fitting error into an optimality guarantee.

The trained neural backend has an explicit computational structure. Under the model's scalar uniform innovation, the expectation of a one-hidden-layer ReLU continuation is analytic. For the scalar investment action, the integrated objective is piecewise quartic, and a finite collection of rationally isolated cubic stationary points supplies exact action-lattice witnesses. This produces a terminating action search for a fitted nonlinear continuation without invoking another unconverged nonconvex optimizer. It is distinct from compiling a native min-plus envelope into an identical neural circuit. The verification theorem nevertheless applies to both neural and conventional candidates.

The policy account is implemented on a locally split, non-tensor observation partition. A one-state residual and value cache encloses the old policy's continuation without a full table of state pairs. A signed sensitivity calculation for the smooth installed reference further retains economically useful continuation direction. We prove its validity and give explicit state-dimension, leaf-count and horizon costs. The original nonlinear investment family has a dimension-uniform smooth-reference stability bound, although the work needed for a chosen approximation accuracy can still grow with dimension. Every later discontinuous acquired actor is handled by the generic cache, not by differentiating away its jumps.

The preceding prospective experiment compares trained ReLU, quadratic and tree-based Bellman continuations in dimensions two, four and eight and horizons two, four and six. The original compiled-witness and tensor-FVI frontends remain in the two-dimensional comparison. Each method pays for its own construction, verification, inference and durable output. Targets require two- or ten-percent expected-cost reductions from installed zero investment. A separately frozen signed-sensitivity extension attains both targets in every declared service; its preceding cache-only failures remain in the record. The complete-policy costs, nonterminal changes and all-in clocks determine what this means economically. They do not establish uniform neural superiority: common verification actions often determine the same returned policy for several generators, and quadratic FVI is an important low-work comparator.

A separate sample-split calculation learns approximately null directions under uncertain action exposure. The guarantee qualifies both the nuisance correction and the remaining Bellman comparison with the same validation confidence box. It therefore concerns inferred structure, rather than merely a hand-selected exact invariant, while retaining the maintained nonlinear drift and innovation model.

The additional comparison in Section~\ref{sec:study57} executes the exact fitted action solver inside the prospective service, rather than assigning old cost observations to new policies. A nested-candidate design keeps the verifier and shared actions fixed, and genuinely different training seeds produce matched fitted objects for the menu and exact-search modes. Theorem~\ref{thm:transfer57} identifies how center optimization loss enters full-action Bellman accuracy; Proposition~\ref{prop:moduli57} supplies explicit action-sensitive acquisition and evaluation moduli. The experiment then separately asks whether the additional witness changes a certified decision and whether its returned policy has a lower actual economic cost.

\paragraph{Related methods.}
Approximate policy iteration and fitted value iteration have established approximation-error analyses \citep{munos2003,munos2008,scherrer2015}. Conservative and safe improvement relate a policy update to an incumbent \citep{kakade2002,pirotta2013,metelli2021}. Tree-based fitted iteration provides a strong conventional comparator \citep{ernst2005}. We do not identify these general principles as neural innovations. Our contribution is the connection among constructive and trained own-future continuations, original feasible witnesses, acquired-state implementation, computable directed verification, and a prospective economic return with complete work accounting. The original action-contrast development also distinguishes residual transport from state bisimulation metrics, reward transformations and value-equivalent model classes \citep{ferns2004,ng1999,grimm2020}.

The article develops this constructive theorem--implementation--economic chain for expected costs. The complete development article and proof companion retain the broader recursive-preference, temporal-self, diffusion and game applications under their original hypotheses. All earlier experiments, including adverse comparisons, are preserved there and in the source archive. Linear action-kernel and sampling identities in the present chain are not silently extended to those different economic evaluations.
